\section{Experimental Setup}

\subsection{Benchmark Scenarios}

The primary evaluation uses seven BSIB scenarios.

\subsubsection{NULL}

\texttt{BSIB-NULL-001} contains no hidden interaction.

It evaluates whether the method avoids reporting interactions when all configurations remain near baseline stochastic risk.

\subsubsection{PAIR-001}

\texttt{BSIB-PAIR-001} contains a strong pairwise interaction between:

\begin{itemize}
\tightlist
\item
  \texttt{tools}; and
\item
  \texttt{structured\_output}.
\end{itemize}

\subsubsection{PAIR-002}

\texttt{BSIB-PAIR-002} contains a weaker pairwise interaction between:

\begin{itemize}
\tightlist
\item
  \texttt{tools}; and
\item
  \texttt{streaming}.
\end{itemize}

This scenario is specifically useful for evaluating stochastic screening stability.

\subsubsection{PAIR-003}

\texttt{BSIB-PAIR-003} contains an interaction between:

\begin{itemize}
\tightlist
\item
  \texttt{parallel\_tools}; and
\item
  \texttt{strict\_schema}.
\end{itemize}

It serves as an additional unseen pairwise scenario.

\subsubsection{OVERLAP}

\texttt{BSIB-OVERLAP-001} contains two pairwise interactions that share one capability:

\begin{itemize}
\tightlist
\item
  \texttt{tools\ +\ streaming}; and
\item
  \texttt{streaming\ +\ strict\_schema}.
\end{itemize}

The scenario tests whether the discovery procedure can recover more than one overlapping pair.

\subsubsection{TRIPLE}

\texttt{BSIB-TRIPLE-001} contains the pure higher-order interaction:

\begin{itemize}
\tightlist
\item
  \texttt{tools\ +\ streaming\ +\ strict\_schema}.
\end{itemize}

Its proper pairs remain near baseline.

The scenario therefore provides no pairwise interaction signal for the hidden triple.

\subsubsection{MIXED}

\texttt{BSIB-MIXED-001} contains both:

\begin{itemize}
\tightlist
\item
  pair: \texttt{tools\ +\ streaming}; and
\item
  triple: \texttt{reasoning\ +\ strict\_schema\ +\ multimodal}.
\end{itemize}

The benchmark is designed to expose masking behavior that affects single-interaction deletion localization.

\subsection{Evaluation Metrics}

The principal metric is exact recovery.

A run is counted as exactly recovered only when the complete set of reported interactions equals the benchmark ground truth.

Additional measures include:

\begin{itemize}
\tightlist
\item
  pairwise recovery;
\item
  higher-order recovery;
\item
  residual-risk classification;
\item
  false-positive runs;
\item
  mean simulator executions;
\item
  median simulator executions;
\item
  minimum and maximum executions;
\item
  95th-percentile execution count;
\item
  reduction relative to exhaustive discovery; and
\item
  wall-clock time in the scaling experiment.
\end{itemize}

\subsection{Unseen Holdout Evaluation}

The primary holdout evaluation uses 1,000 unseen seeds per benchmark scenario.

The holdout seed range is:

\[
31001,\ldots,32000.
\]

Seven scenarios are evaluated, producing a total of

\[
7 \times 1000 = 7000
\]

SCIF v0.0.4 holdout runs.

Parameters were not modified after observing the holdout results.

This is important because one rare weak-pair miss was retained and analyzed rather than being used to retune the algorithm.

\subsection{Baseline Comparison}

Four discovery methods are compared:

\begin{enumerate}
\def\labelenumi{\arabic{enumi}.}
\tightlist
\item
  SCIF v0.0.3;
\item
  deletion-based localization;
\item
  exhaustive discovery; and
\item
  SCIF v0.0.4.
\end{enumerate}

The comparison study uses 100 seeds from:

\[
41001,\ldots,41100.
\]

The selected scenarios are:

\begin{itemize}
\tightlist
\item
  PAIR-002;
\item
  OVERLAP;
\item
  TRIPLE; and
\item
  MIXED.
\end{itemize}

For pairwise scenarios, exhaustive discovery evaluates configurations through order two.

For scenarios containing a triple, exhaustive discovery evaluates configurations through order three.

With eight capabilities this corresponds to:

\[
1 + 8 + \binom{8}{2}
=
37
\]

configurations for exhaustive order-two discovery, or 37,000 simulator executions at 1,000 trials per configuration.

For exhaustive order-three discovery:

\[
1 + 8 + \binom{8}{2} + \binom{8}{3}
=
93
\]

configurations are evaluated, corresponding to 93,000 simulator executions.

For adaptive methods, execution cost is the actual number of simulator invocations performed by the method, including screening, retesting, confirmation, residual probing, and localization where applicable. The exhaustive baselines instead use the fixed 1,000-trial budget for every enumerated configuration.

\subsection{Ablation Study}

The ablation study uses 100 seeds:

\[
51001,\ldots,51100.
\]

Five representative scenarios are evaluated:

\begin{itemize}
\tightlist
\item
  NULL;
\item
  PAIR-002;
\item
  OVERLAP;
\item
  TRIPLE; and
\item
  MIXED.
\end{itemize}

Three algorithmic stages are compared:

\subsubsection{Stage A --- V3 Pairwise Only}

Only confirmed pairwise discoveries are considered.

\subsubsection{Stage B --- V3 + Residual Detection}

Pairwise discovery is followed by residual-risk classification.

This stage determines whether higher-order escalation is needed but does not yet localize a higher-order interaction.

\subsubsection{Stage C --- Full V4}

The complete pipeline includes residual higher-order localization.

The ablation records:

\begin{itemize}
\tightlist
\item
  stage success;
\item
  false residual escalations;
\item
  missed escalations;
\item
  false higher-order candidates; and
\item
  mean execution cost.
\end{itemize}

\subsection{Sensitivity Study}

The sensitivity study uses 100 seeds:

\[
61001,\ldots,61100.
\]

Four scenarios are tested:

\begin{itemize}
\tightlist
\item
  NULL;
\item
  PAIR-002;
\item
  TRIPLE; and
\item
  MIXED.
\end{itemize}

The frozen baseline configuration is:

\begin{itemize}
\tightlist
\item
  \texttt{initial\_trials\ =\ 100};
\item
  \texttt{min\_residual\_increment\ =\ 0.10}; and
\item
  \texttt{higher\_order\_min\_removal\_drop\ =\ 0.10}.
\end{itemize}

The evaluated alternatives are:

\begin{itemize}
\tightlist
\item
  initial trials = 50;
\item
  initial trials = 200;
\item
  residual increment = 0.05;
\item
  residual increment = 0.15;
\item
  removal drop = 0.05; and
\item
  removal drop = 0.15.
\end{itemize}

The purpose is sensitivity characterization rather than post-holdout parameter optimization.

\subsection{Scaling Study}

Scaling is evaluated using benchmark variants with:

\begin{itemize}
\tightlist
\item
  8 capabilities;
\item
  12 capabilities;
\item
  16 capabilities; and
\item
  20 capabilities.
\end{itemize}

Each benchmark retains:

\begin{itemize}
\tightlist
\item
  one pairwise interaction; and
\item
  one three-way interaction.
\end{itemize}

Additional capabilities are non-interacting noise capabilities.

Twenty seeds are evaluated for each size:

\[
71001,\ldots,71020.
\]

For \(n\) capabilities, exhaustive order-three discovery evaluates

\[
1+n+\binom{n}{2}+\binom{n}{3}
\]

configurations.

At 1,000 trials per configuration, exhaustive execution counts are:

\begin{itemize}
\tightlist
\item
  93,000 for 8 capabilities;
\item
  299,000 for 12 capabilities;
\item
  697,000 for 16 capabilities; and
\item
  1,351,000 for 20 capabilities.
\end{itemize}

\subsection{Reproducibility}

The reproducibility package contains benchmark definitions, simulator implementation, discovery algorithms, automated tests, experiment scripts, aggregated result files, and publication figure-generation scripts.

All Python development commands are executed through a locked \texttt{uv} environment. Repository-identifying information is omitted during anonymous review.
